%--- ХУРААНГУЙ ---
\addchaptertocentry{Хураангуй}

\begin{center}
{\scshape\Large \univname\par}
{\scshape\large \facname\par}\vspace{0.5cm}
{\huge\textbf{Хураангуй} \par}
\bigskip
{\Large{\ttitle} \par}
\bigskip
{\normalsize \shortname \par}
\mailaddress
\end{center}

\bigskip

QR код ашиглан хоол захиалах систем нь ресторан, кафе болон хоолны газруудын захиалга
өгөх процессийг автоматжуулах зорилготой. Уламжлалт аргаар зөөгчөөр дамжуулан захиалга
авдаг схем нь мэдээлэл дамжуулахад хүний оролцоог шаарддаг тул удаашрал, алдааны
эрсдэл дагуулдаг. Энэхүү судалгааны ажилд QR кодоор идэвхждэг мобайл захиалгын системийг
загварчилж, React Native мобайл аппликейшн, Node.js/Express RESTful API сервер, PostgreSQL
өгөгдлийн сан болон Socket.IO WebSocket технологиудын уялдаа бүхий 4 давхаргат архитектурт
хэрэгжүүллээ.

Хэрэглэгч ширээн дэлгэц дээрх QR кодыг смартфоноор уншуулахад тухайн рестораны
дижитал цэс нэн даруй нээгдэж, хоол сонгон захиалга өгөхөд захиалга нь WebSocket
протоколоор ажилтны дэлгэцэнд бодит цаг хугацаанд дамжина. Улаанбаатар хотын 3 рестораны
1,240 захиалгын мэдээлэлд тулгуурласан туршилтаар захиалгын нийт хугацаа 83\%,
захиалгын алдааны түвшин 91\% тус тус буурч, 60 зочин болон 18 ажилтны нийт сэтгэл
ханамжийн үнэлгээ 4.5/5.0 оноонд хүрсэн болно.

\bigskip
\noindent\textit{\textbf{Түлхүүр үгс:} QR код, мобайл захиалга, рестораны мэдээллийн систем, WebSocket, бодит цагийн мэдэгдэл, React Native}

\bigskip

\begin{center}
{\huge\textbf{Abstract} \par}
\bigskip
{\Large{\ttitleng} \par}
\bigskip
{\normalsize \shortname \par}
\end{center}

\bigskip

This research presents the design, implementation, and evaluation of a QR code-based mobile
table ordering system for restaurants. The proposed system eliminates the traditional waiter-mediated 
ordering process by enabling customers to scan a QR code at their table, browse a digital menu, 
and place orders that are instantly delivered to staff via WebSocket protocol in real time. The system 
integrates a React Native mobile client, Node.js/Express RESTful API, PostgreSQL database, and 
Socket.IO WebSocket server in a four-layer architecture.

Experimental results from 1,240 orders across three Ulaanbaatar restaurants demonstrate an 83\% 
reduction in total ordering time (346 to 59 seconds), a 91\% decrease in order error rate (23\% to 
2\%), and a total satisfaction score of 4.5/5.0 from 60 customers and 18 staff members. The solution 
is validated as a practical and scalable deployment option for the Mongolian restaurant industry.

\bigskip
\noindent\textit{\textbf{Keywords:} QR code, mobile ordering, restaurant information system, WebSocket, real-time notification, React Native}

\clearpage
